
\documentclass[11pts,a4paper,amsmath,amssymb,floatfix]{book}

%\usepackage[gather]{chapterbib}

\usepackage{graphicx,wrapfig}% Include figure files
%\usepackage{dcolumn,enumerate}% Align table columns on decimal point
\usepackage{enumerate}% Align table columns on decimal point
\usepackage{bm,dpfloat}% bold math
\usepackage[pdftex,bookmarks,colorlinks=true,urlcolor=rltblue,citecolor=blue]{hyperref}
\usepackage{amsfonts,amsmath,amsthm,amssymb,stmaryrd,indentfirst}
\usepackage{times,psfrag,pdfpages}
%\usepackage{natbib} 
\usepackage{color}
\usepackage{units}
\usepackage{rotating}
\usepackage{multirow}

\usepackage{pifont}
\usepackage{subfigure}
\usepackage{subeqnarray}
\usepackage{ifthen}

\usepackage{supertabular}
\usepackage{moreverb}
\usepackage{fancyvrb}
\usepackage{listings}
\usepackage{palatino}
%\usepackage{doi}
\usepackage{longtable}
\usepackage{float}
\usepackage{perpage}
\MakeSorted{figure}
%\usepackage{pdflscape}

% Appendix allows the inclusion of a line for "appendices", which otherwise won't get the right page number. You must
% use the [toc] option.
%\usepackage[toc]{appendix}% Text + Eqn box:

\usepackage[framemethod=TikZ]{./mdframed/mdframed}
\usepackage{lipsum}
\mdfdefinestyle{JFrame}{%
    linecolor=blue,
    outerlinewidth=1pt,
    roundcorner=20pt,
    innertopmargin=\baselineskip,
    innerbottommargin=\baselineskip,
    innerrightmargin=20pt,
    innerleftmargin=20pt, 
    backgroundcolor=blue!10!white}

%\usepackage[framemethod=default]{./mdframed/mdframed}
%\usepackage{showexpl}
%\mdfdefinestyle{exampledefault}{%
%rightline=true,innerleftmargin=10,innerrightmargin=10,
%frametitlerule=true,frametitlerulecolor=green,
%frametitlebackgroundcolor=yellow,
%frametitlerulewidth=2pt}

\definecolor{rltblue}{rgb}{0,0,0.75}

%\usepackage{natbib}
\usepackage{fancyhdr} %%%%
\pagestyle{fancy}%%%%
% with this we ensure that the chapter and section
% headings are in lowercase
%%%%\renewcommand{\chaptermark}[1]{\markboth{#1}{}}
\renewcommand{\sectionmark}[1]{\markright{\thesection\ #1}}
\fancyhf{} %delete the current section for header and footer
%\fancyhead[LE,RO]{\bfseries\thepage}
%\fancyfoot[LE,RO]{\bfseries\thepage}
\fancyhead[LO]{\bfseries\rightmark}
\fancyhead[RE]{\bfseries\leftmark}
\renewcommand{\headrulewidth}{0.5pt}
% make space for the rule
\fancypagestyle{plain}{%
\fancyhead{} %get rid of the headers on plain pages
\renewcommand{\headrulewidth}{0pt} % and the line
}

\def\newblock{\hskip .11em plus .33em minus .07em}
\usepackage{color}

\theoremstyle{definition}
\newtheorem{exmp}{Example}[chapter]

% For theorems
\newtheorem{theorem}{Theorem}[section]
\newtheorem{corollary}{Corollary}[theorem]
\newtheorem{lemma}[theorem]{Lemma}

\usepackage{makeidx}
\makeindex

\setlength\textwidth      {16.cm}
\setlength\textheight     {23.0cm}
\setlength\oddsidemargin  {-0.3cm}
\setlength\evensidemargin {0.3cm}

\setlength\headheight{14.49998pt} 
\setlength\topmargin{0.0cm}
\setlength\headsep{1.cm}
\setlength\footskip{-0.cm}
\setlength\parskip{0pt}

%%%
%%% Headers and Footers
%\lhead[\text{\small{WORK IN PROGRESS}}] {\text{\small{ICON/AMCG}}} 
%-\lhead[] {\text{\small{ICON/AMCG}}} 
%-\chead[\text{\small{FETCH Manual}}]  {\text{\small{FETCH Manual}}} %{\text{\small{JAERI Analysis Report Issue}}}
%-\rhead[\text{\small{Commercial-in Confidence}}]{}
\cfoot[\thepage]{\thepage}
%-\rfoot[]{}%{\text{\small{\today}}}
%\cfoot[\text{\small{\today}}]{}% {\text{\small{\today}}}
%\cfoot[\text{\small{January 2006}}] {\text{\small{January 2006}}}
%-\lfoot []{}%{\text{\small{Numerical Analysis}}}
%-\renewcommand{\headrulewidth}{0.8pt}

%%%
%%% space between lines
%%%
\renewcommand{\baselinestretch}{1.5}

\newenvironment{VarDescription}[1]%
  {\begin{list}{}{\renewcommand{\makelabel}[1]{\textbf{##1:}\hfil}%
    \settowidth{\labelwidth}{\textbf{#1:}}%
    \setlength{\leftmargin}{\labelwidth}\addtolength{\leftmargin}{\labelsep}}}%
  {\end{list}}

%%%%%%%%%%%%%%%%%%%%%%%%%%%%%%%%%%%%%%%%%%%
%%%%%%                              %%%%%%%
%%%%%%      NOTATION SECTION        %%%%%%%
%%%%%%                              %%%%%%%
%%%%%%%%%%%%%%%%%%%%%%%%%%%%%%%%%%%%%%%%%%%

% Text abbreviations.
\newcommand{\ie}{{\em{i.e., }}}
\newcommand{\eg}{{\em{e.g., }}}
\newcommand{\cf}{{\em{cf., }}}
\newcommand{\fetch}{FETCH }
\newcommand{\fmips}{FETCH-MIPS }
\newcommand{\bw}{B$\$$W }
\newcommand{\fluidity}{Fluidity }
\newcommand{\event}{{EVENT }}
\newcommand{\gem}{GEM }
\newcommand{\radiant}{RADIANT }
\newcommand{\wrt}{with respect to}
\newcommand{\lhs}{left hand side}
\newcommand{\rhs}{right hand side}
\newcommand{\vv}{V$\&$V }
\newcommand{\cse}{CS$\&$E }
\newcommand{\amcg}{\href{http://amcg.ese.ic.ac.uk}{AMCG }}
% Commands definining mathematical notation.

% This is for quantities which are physically vectors.
\renewcommand{\vec}[1]{{\mbox{\boldmath$#1$}}}
% Physical rank 2 tensors
\newcommand{\tensor}[1]{\overline{\overline{#1}}}
% This is for vectors formed of the value of a quantity at each node.
\newcommand{\dvec}[1]{\underline{#1}}
% This is for matrices in the discrete system.
\newcommand{\mat}[1]{\mathrm{#1}}


\DeclareMathOperator{\sgn}{sgn}

%\newcommand\qed{\hfill\mbox{$\Box$}}
\newcommand{\re}{{\mathrm{I}\hspace{-0.2em}\mathrm{R}}}
\newcommand{\inner}[2]{\langle#1,#2\rangle}
\renewcommand\leq{\leqslant}
\renewcommand\geq{\geqslant}
\renewcommand\le{\leqslant}
\renewcommand\ge{\geqslant}
\renewcommand\epsilon{\varepsilon}
\newcommand\eps{\varepsilon}
\newcommand{\bmF}{\vec{F}}
\newcommand{\bmphi}{\vec{\phi}}
\newcommand{\bmn}{\vec{n}}
\newcommand{\bmns}{{\textrm{\scriptsize{\boldmath $n$}}}}
\newcommand{\bmi}{\vec{i}}
\newcommand{\bmj}{\vec{j}}
\newcommand{\bmk}{\vec{k}}
\newcommand{\bmx}{\vec{x}}
\newcommand{\bmu}{\vec{u}}
\newcommand{\bmv}{\vec{v}}
\newcommand{\bmr}{\vec{r}}
\newcommand{\bma}{\vec{a}}
\newcommand{\bmg}{\vec{g}}
\newcommand{\bmU}{\vec{U}}
\newcommand{\bmI}{\vec{I}}
\newcommand{\bmq}{\vec{q}}
\newcommand{\bmT}{\vec{T}}
\newcommand{\bmM}{\vec{M}}
\newcommand{\bmtau}{\vec{\tau}}
\newcommand{\bmOmega}{\vec{\Omega}}
\newcommand{\pp}{\partial}
\newcommand{\kaptens}{\tensor{\kappa}}
\newcommand{\tautens}{\tensor{\tau}}
\newcommand{\sigtens}{\tensor{\sigma}}
\newcommand{\etens}{\tensor{\dot\epsilon}}
\newcommand{\ktens}{\tensor{k}}
\newcommand{\half}{{\textstyle \frac{1}{2}}}
\newcommand{\tote}{E}
\newcommand{\inte}{e}
\newcommand{\strt}{\dot\epsilon}
\newcommand{\modu}{|\bmu|}
\newcommand{\summation}{\sum\limits}
% Derivatives
\renewcommand{\d}{\mathrm{d}}
\newcommand{\D}{\mathrm{D}}
\newcommand{\ddx}[2][x]{\frac{\d#2}{\d#1}}
\newcommand{\ddxx}[2][x]{\frac{\d^2#2}{\d#1^2}}
\newcommand{\ddt}[2][t]{\frac{\d#2}{\d#1}}
\newcommand{\ddtt}[2][t]{\frac{\d^2#2}{\d#1^2}}
\newcommand{\ppx}[2][x]{\frac{\partial#2}{\partial#1}}
\newcommand{\ppxx}[2][x]{\frac{\partial^2#2}{\partial#1^2}}
\newcommand{\ppt}[2][t]{\frac{\partial#2}{\partial#1}}
\newcommand{\pptt}[2][t]{\frac{\partial^2#2}{\partial#1^2}}
\newcommand{\DDx}[2][x]{\frac{\D#2}{\D#1}}
\newcommand{\DDxx}[2][x]{\frac{\D^2#2}{\D#1^2}}
\newcommand{\DDt}[2][t]{\frac{\D#2}{\D#1}}
\newcommand{\DDtt}[2][t]{\frac{\D^2#2}{\D#1^2}}
% Norms
\newcommand{\Ltwo}{\ensuremath{L_2} }
% Basis functions
\newcommand{\Qone}{\ensuremath{Q_1} }
\newcommand{\Qtwo}{\ensuremath{Q_2} }
\newcommand{\Qthree}{\ensuremath{Q_3} }
\newcommand{\QN}{\ensuremath{Q_N} }
\newcommand{\Pzero}{\ensuremath{P_0} }
\newcommand{\Pone}{\ensuremath{P_1} }
\newcommand{\Ptwo}{\ensuremath{P_2} }
\newcommand{\Pthree}{\ensuremath{P_3} }
\newcommand{\PN}{\ensuremath{P_N} }
\newcommand{\Poo}{\ensuremath{P_1P_1} }
\newcommand{\PoDGPt}{\ensuremath{P_{-1}P_2} }

\newcommand{\metric}{\tensor{M}}
\newcommand{\configureflag}[1]{\texttt{#1}}

% Units
\newcommand{\m}[1][]{\unit[#1]{m}}
\newcommand{\km}[1][]{\unit[#1]{km}}
\newcommand{\s}[1][]{\unit[#1]{s}}
\newcommand{\invs}[1][]{\unit[#1]{s}\ensuremath{^{-1}}}
\newcommand{\ms}[1][]{\unit[#1]{m\ensuremath{\,}s\ensuremath{^{-1}}}}
\newcommand{\mss}[1][]{\unit[#1]{m\ensuremath{\,}s\ensuremath{^{-2}}}}
\newcommand{\K}[1][]{\unit[#1]{K}}
\newcommand{\PSU}[1][]{\unit[#1]{PSU}}
\newcommand{\Pa}[1][]{\unit[#1]{Pa}}
\newcommand{\kg}[1][]{\unit[#1]{kg}}
\newcommand{\rads}[1][]{\unit[#1]{rad\ensuremath{\,}s\ensuremath{^{-1}}}}
\newcommand{\kgmm}[1][]{\unit[#1]{kg\ensuremath{\,}m\ensuremath{^{-2}}}}
\newcommand{\kgmmm}[1][]{\unit[#1]{kg\ensuremath{\,}m\ensuremath{^{-3}}}}
\newcommand{\Nmm}[1][]{\unit[#1]{N\ensuremath{\,}m\ensuremath{^{-2}}}}

% Dimensionless numbers
\newcommand{\dimensionless}[1]{\mathrm{#1}}
\renewcommand{\Re}{\dimensionless{Re}}
\newcommand{\Ro}{\dimensionless{Ro}}
\newcommand{\Fr}{\dimensionless{Fr}}
\newcommand{\Bu}{\dimensionless{Bu}}
\newcommand{\Ri}{\dimensionless{Ri}}
\renewcommand{\Pr}{\dimensionless{Pr}}
\newcommand{\Pe}{\dimensionless{Pe}}
\newcommand{\Ek}{\dimensionless{Ek}}
\newcommand{\Gr}{\dimensionless{Gr}}
\newcommand{\Ra}{\dimensionless{Ra}}
\newcommand{\Sh}{\dimensionless{Sh}}
\newcommand{\Sc}{\dimensionless{Sc}}


% Journals
\newcommand{\IJHMT}{{\it International Journal of Heat and Mass Transfer}}
\newcommand{\NED}{{\it Nuclear Engineering and Design}}
\newcommand{\ICHMT}{{\it International Communications in Heat and Mass Transfer}}
\newcommand{\NET}{{\it Nuclear Engineering and Technology}}
\newcommand{\HT}{{\it Heat Transfer}}   
\newcommand{\IJHT}{{\it International Journal for Heat Transfer}}


\newcommand{\frc}{\displaystyle\frac}
\newcommand{\red}{\textcolor{red}}
\newcommand{\blue}{\textcolor{blue}}
\newcommand{\green}{\textcolor{green}}
\newcommand{\purple}{\textcolor{purple}}
\newcommand{\PNDG}[3][error]{P$_{#1}$DG-P$_{#2}$DG$_{#3}$}
\newcommand{\Partial}[3][error]{\left(\frc{\partial #1}{\partial #2}\right)_{#3}}
%{P$_{#1}$DG-P$_{#2}$DG}%$_{#3}$}

%%%%%%%%%%%%%%%%%%%%%%%%%%%%%%%%%%%%%%%%%%%
%%%%%%                              %%%%%%%
%%%%%% END OF THE NOTATION SECTION  %%%%%%%
%%%%%%                              %%%%%%%
%%%%%%%%%%%%%%%%%%%%%%%%%%%%%%%%%%%%%%%%%%%

% Cause numbering of subsubsections. 
%\setcounter{secnumdepth}{8}
%\setcounter{tocdepth}{8}

\setcounter{secnumdepth}{4}%
\setcounter{tocdepth}{4}%

\newcounter{qcounter}
\includeonly{Module_01, Module_02}
%\includeonly{model_introduction, model_reservoir_equations, model_description}
\DeclareMathAlphabet{\mathpzc}{OT1}{pzc}{m}{it}

\begin{document}

\vspace{4cm}

\begin{titlepage}
\begin{center}
\bigskip

   \includegraphics[width=15cm,clip]{./FigBanner/UoAHorizBanner}

\bigskip
\vspace{3.5cm}

   {\bf{\Huge Integrated Multi-Fluid Flows Transport and Thermodynamic Models}}

\bigskip
   {\bf{\huge Model Documentation}}

\vspace{3.5cm}

\end{center}

\begin{tabular}{l c l}
{\bf{\large Author:}}                          &     & J.L.M.A. Gomes \\
                                               &     & Fluid and Structure Group \\
                                               &     & School of Engineering \\
                                               &     & University of Aberdeen \\
\bigskip
{\bf{\large Date:}}                            &     &   \today\\
                                               &     &         \\
{\bf{\large Revision 1:}}                      &     &   \today\\
                                               &     &         \\
\end{tabular}
\vspace{3.5cm}

\bigskip

\bigskip




\end{titlepage}

\setcounter{page}{1}

\tableofcontents
\vfill

\pagebreak
\listoftables
\vfill
\pagebreak
\listoffigures
\vfill
\pagebreak


%%%
%%% CHAPTER
%%%
\chapter{Thermodynamic Formulation}

 
%%% SECTION
\section{Mass Balance}
For a system with $n_{p}$ phases and $n_{c}$ components, the mass of each $j$ phase is,
\begin{equation}
m^{j} = \summation_{i=1}^{n_{c}}m_{i}^{j}\;\;\;\text{ for }j=1,2,\cdots,n_{p}
\label{Eqn_MassBalance_Phase}
\end{equation}

And the mass of each component $i\left(=1,2,\cdots,n_{c}\right)$,
\begin{equation}
m_{i} = m_{i}^{1} + m_{i}^{2} + \cdots + m_{i}^{n_{p}} = \summation_{j=1}^{n_{p}}m_{i}^{j}
\label{Eqn_MassBalance_Component}
\end{equation}

The total mass $m$ can be define as,
\begin{eqnarray}
m & = & m^{1} + m^{2} + \cdots + m^{n_{p}} = \summation_{j=1}^{n_{p}}m^{j} \nonumber \\
  & = & m_{1} + m_{2} + \cdots + m_{n_{c}} = \summation_{i=1}^{n_{c}}m_{i} \nonumber 
\label{Eqn_MassBalance_Mass}
\end{eqnarray}

Defining the intensive properties:
\begin{eqnarray}
&& w_{i}^{j} = \frc{m_{i}^{j}}{m^{j}}, \hspace{2cm} z_{i} = \frc{m_{i}}{m},\label{Eqn_MassBalance_MassFraction} \\
&& \Pi^{j} = \frc{m^{j}}{m}\label{Eqn_MassBalance_PhaseMassFraction} 
\end{eqnarray}
where $w_{i}^{j}$ is the mass fraction of component $i$ in phase $j$, $\Pi^{j}$ is the mass fraction of phase $j$ and $z_{i}$ is the overall feed mass fraction of component $i$. From Eqns.~\ref{Eqn_MassBalance_Mass}-~\ref{Eqn_MassBalance_PhaseMassFraction}, we can define normalised quantities, i.e., fractions, and we can add the following constraints:
\begin{eqnarray}
&&w_{n_{c}}^{j} = 1 - \summation_{i=1}^{n_{c}-1}w_{i}^{j},\hspace{1cm}  z_{n_{c}} = 1 - \summation_{i=1}^{n_{c}-1}z_{i},  \label{Eqn_MassBalance_MassFraction2} \\
&&\Pi^{n_{p}} = 1 - \left(\Pi^{1} + \Pi^{2} + \cdots + \Pi^{n_{p}-1}\right) = 1 - \summation_{j=1}^{n_{p}-1} \Pi^{j}\hspace{1cm} j=1,2,\cdots,n_{p} \label{Eqn_MassBalance_PhaseMassFraction2}
\end{eqnarray} 
with
\begin{displaymath}
  z_{i} = \frc{m_{i}}{m} = \frc{\summation_{j=1}^{n_{p}}m_{i}^{j}}{m} \frc{m^{j}}{m^{j}} = \underbrace{\frc{m^{j}}{m}}_{\red{\Pi^{j}}} \cdot \overbrace{\frc{\summation_{j=1}^{n_{p}}m_{i}^{j}}{m^{j}}}^{\red{\summation_{j=1}^{n_{p}}w_{i}^{j}}} 
\end{displaymath}
leading to
\begin{equation}
z_{i} = \Pi^{1}w_{i}^{1} + \Pi^{2}w_{i}^{2} + \cdots + \Pi^{n_{p}}w_{i}^{n_{p}}
\label{Eqn_MassBalance_FeedMassFractionConstraint}
\end{equation}
with
\begin{equation}
0\leq\Pi^{j}\leq 1 \hspace{3cm} 0\leq w_{i}^{j}\leq 1
\end{equation}

If the solution is contained within \red{$n_{p}$} phases, the inequality can be rewritten as,
\begin{equation} 
0 < \Pi^{j} < 1
\end{equation}
Thus
\begin{displaymath}
\Pi^{k} = 1 - \summation_{j=1,j\neq k}^{n_{p}} \Pi^{j} \neq 0
\end{displaymath}
Therefore from Eqn.~\ref{Eqn_MassBalance_FeedMassFractionConstraint},
\begin{equation}
w_{i}^{k} = \frc{z_{i}-\summation_{j=1,j\neq k}^{n_{p}}\Pi^{j}w_{i}^{j}}{\Pi^{k}} = \frc{z_{i} - \summation_{j=1,j\neq k}^{n_{p}}\Pi^{j}w_{i}^{j}}{1-\summation_{j=1,j\neq k}^{n_{p}}\Pi^{j}}
\label{Eqn_MassBalance_FeedMassFractionConstraint2}
\end{equation}


\begin{mdframed}style=exampledefault,frametitle={Inhomogeneous linear}]%[style=JFrame]
For 2 phases, $\Pi^{1}=L$ and $\Pi^{2}=V$, Eqn.~\ref{Eqn_MassBalance_FeedMassFractionConstraint} becomes:
\begin{displaymath}
w_{i}^{V} = \frc{z_{i}-Lw_{i}^{L}}{1-L}
\end{displaymath}
for $i=1,2,\cdots,n_{c}$, and for 3 phases,$\Pi^{1}=L$, $\Pi^{2}=V$ and $\Pi^{3}=H$:
\begin{displaymath}
w_{i}^{H} = \frc{z_{i}-\left(V w_{i}^{V} + L w_{i}^{L}\right)}{1 - \left(V w_{i}^{V} + L w_{i}^{L}\right)}
\end{displaymath}
\end{mdframed}


%

%%%
%%% CHAPTER
%%%
\section{Model Formulation}
 
%%% SUB-SECTION
\subsection{Mass Balance}
For a system with $n_{p}$ phases and $n_{c}$ components, the mass of each phase $j$ is,
\begin{equation}
m^{j} = \summation_{i=1}^{n_{c}}m_{i}^{j}\;\;\;\text{ for }j=1,2,\cdots,n_{p}
\label{Eqn_MassBalance_Phase}
\end{equation}

And the mass of each component $i\left(=1,2,\cdots,n_{c}\right)$,
\begin{equation}
m_{i} = m_{i}^{1} + m_{i}^{2} + \cdots + m_{i}^{n_{p}} = \summation_{j=1}^{n_{p}}m_{i}^{j}
\label{Eqn_MassBalance_Component}
\end{equation}

The total mass $m$ can be define as,
\begin{equation}
m - m^{1} + m^{2} + \cdots + m^{n_{p}} = \summation_{j=1}^{n_{p}}m^{j} = \summation_{i=1}^{n_{c}}m_{i}
\label{Eqn_MassBalance_Mass}
\end{equation}

Defining the intensive properties:
\begin{eqnarray}
w_{i}^{j} = \frc{m_{i}^{j}}{m^{j}} \label{Eqn_MassBalance_MassFraction} \\
z_{i} = \frc{m_{i}}{m} \label{Eqn_MassBalance_MassFractionFeeding} \\
\Pi^{j} = \frc{m^{j}}{m}\label{Eqn_MassBalance_PhaseMassFraction} 
\end{eqnarray}
where $w_{i}^{j}$ is the mass fraction of component $i$ in phase $j$, $\Pi^{j}$ is the mass fraction of phase $j$ and $z_{i}$ is the overall feed mass fraction of component $i$. From Eqns.~\ref{Eqn_MassBalance_Mass}-~\ref{Eqn_MassBalance_PhaseMassFraction}, we can define normalised quantities, i.e., fractions, and we can add the following constraints:
\begin{eqnarray}
&&w_{n_{c}}^{j} = 1 - \summation_{i=1}^{n_{c}-1}w_{i}^{j}  \hspace{3cm} j=1,2,\cdots,n_{p}  \label{Eqn_MassBalance_MassFraction2} \\
&&z_{n_{c}} = 1 - \summation_{i=1}^{n_{c}-1}z_{1}  \\
&&\Pi^{n_{p}} = 1 - \left(\Pi^{1} + \Pi^{2} + \cdots + \Pi^{n_{p}-1}\right) = 1 - \summation_{j=1}^{n_{p}-1} \Pi^{j}\label{Eqn_MassBalance_PhaseMassFraction2}
\end{eqnarray} 
with
\begin{displaymath}
  z_{i} = \frc{m_{i}}{m} = \frc{\summation_{j=1}^{n_{p}}m_{i}^{j}}{m} \frc{m^{j}}{m^{j}} = \underbrace{\frc{m^{j}}{m}}_{\red{\Pi^{j}}} \cdot \overbrace{\frc{\summation_{j=1}^{n_{p}}m_{i}^{j}}{m^{j}}}^{\red{\summation_{j=1}^{n_{p}}w_{i}^{j}}} 
\end{displaymath}
leading to
\begin{equation}
z_{i} = \Pi^{1}w_{i}^{1} + \Pi^{2}w_{i}^{2} + \cdots + \Pi^{n_{p}}w_{i}^{n_{p}}
\label{Eqn_MassBalance_FeedMassFractionConstraint}
\end{equation}
with
\begin{equation}
0\leq\Pi^{j}\leq 1 \hspace{3cm} 0\leq w_{i}^{j}\leq 1
\end{equation}

If the solution is contained within \red{$n_{p}$} phases, the inequality can be rewritten as,
\begin{equation}
0 < \Pi^{j} < 1
\end{equation}
Thus 
\begin{displaymath}
\Pi^{k} = 1 - \summation_{j=1,j\neq k}^{n_{p}} \Pi^{j} \neq 0
\end{displaymath}
Therefore from Eqn.~\ref{Eqn_MassBalance_FeedMassFractionConstraint},
\begin{equation}
w_{i}^{k} = \frc{z_{i}+\summation_{j=1,j\neq k}^{n_{p}}\Pi^{j}w_{i}^{j}}{\Pi^{k}} = \frc{z_{i} + \summation_{j=1,j\neq k}^{n_{p}}\Pi^{j}w_{i}^{j}}{1-\summation_{j=1,j\neq k}^{n_{p}}\Pi^{j}}
\label{Eqn_MassBalance_FeedMassFractionConstraint2}
\end{equation}



%\begin{mdframed}[style=JFrame]
For 2 phases, $\Pi^{1}=L$ and $\Pi^{2}=V$, Eqn.~\ref{Eqn_MassBalance_FeedMassFractionConstraint} becomes:
\begin{displaymath}
w_{i}^{V} = \frc{z_{i}-Lw_{i}^{L}}{1-L}
\end{displaymath}
for $i=1,2,\cdots,n_{c}$, and for 3 phases,$\Pi^{1}=L$, $\Pi^{2}=V$ and $\Pi^{3}=H$:
\begin{displaymath}
w_{i}^{H} = \frc{z_{i}-\left(V w_{i}^{V} + L w_{i}^{L}\right)}{1 - \left(V w_{i}^{V} + L w_{i}^{L}\right)}
\end{displaymath}
%\end{mdframed}

%% SUB-SECTION
\subsection{Free Gibbs Energy}
For equilibrium problems (e.g., vapour-liquid equilibrium, VLE), the total free Gibbs energy can be expressed as,
\begin{equation}
G = \summation_{j=1}^{n_{p}}\left(\summation_{i=1}^{n_{c}} m_{i}^{j}\mu_{i}^{j}\right)
\label{MassGibbs_Definition}
\end{equation} 
this represents
\begin{displaymath}
G = \summation_{i=1}^{n_{c}} m_{i}^{1}\mu_{i}^{1} + \summation_{i=1}^{n_{c}} m_{i}^{2}\mu_{i}^{2} + \cdots + \summation_{i=1}^{n_{c}} m_{i}^{n_{p}}\mu_{i}^{n_{p}}
\end{displaymath}
The chemical potential $\left(\mu\right)$ can be written as a function of the mass of each component at each phase,
\begin{eqnarray}
\mu_{i}^{j} &=& \mu_{i}^{j}\left(m_{1}^{j}, m_{2}^{j}, \cdots, m_{n_{c}}^{j}\right) = \mu_{i}^{j}\left(\frc{m_{1}^{j}}{m^{j}}, \frc{m_{2}^{j}}{m^{j}}, \cdots, \frc{m_{n_{c}}^{j}}{m^{j}}\right)\nonumber \\
 &=& \mu_{i}^{j}\left(w_{1}^{j}, w_{2}^{j}, \cdots, w_{n_{c}}^{j}\right)\label{chempotential_functional}
\end{eqnarray}
Dividing Eqn.~\ref{MassGibbs_Definition} by $m$,
\begin{equation}
g = \frc{G}{m} = \summation_{j=1}^{n_{p}}\left(\summation_{i=1}^{n_{c}}\frc{m_{i}^{j}}{m}\mu_{i}^{j}\right) = \summation_{i=1}^{n_{c}}\frc{m_{i}^{1}}{m^{1}} \cdot \frc{m^{1}}{m} \mu_{i}^{1} + \cdots + \summation_{i=1}^{n_{c}}\frc{m_{i}^{n_{p}}}{m^{n_{p}}} \cdot \frc{m^{n_{p}}}{m} \mu_{i}^{n_{p}}
\label{MassGibbs_Definition2}
\end{equation}
And we rewrite Eqn.~\ref{MassGibbs_Definition2} as,
\begin{equation}
g = \summation_{i=1}^{n_{c}}\Pi^{1}w_{i}^{1}\mu_{i}^{1} + \summation_{i=1}^{n_{c}}\Pi^{2}w_{i}^{2}\mu_{i}^{2} + \cdots + \summation_{i=1}^{n_{c}}\Pi^{n_{p}}w_{i}^{n_{p}}\mu_{i}^{n_{p}} = \summation_{j=1}^{n_{p}}\left[\summation_{i}^{n_{c}}\Pi^{j}w_{i}^{j}\mu_{i}^{j}\right] \label{MassGibbs_Definition3}
\end{equation}


%\begin{mdframed}[style=JFrame]
For 2 ($L$ and $V$) phases,
\begin{displaymath}
g = \summation_{i=1}^{n_{c}}Lw_{i}^{L}\mu_{i}^{L} + \summation_{i=1}^{n_{c}}Vw_{i}^{V}\mu_{i}^{V}
\end{displaymath}
and 3 ($L$, $V$ and $H$) phases,
\begin{displaymath}
g = \summation_{i=1}^{n_{c}}Lw_{i}^{L}\mu_{i}^{L} + \summation_{i=1}^{n_{c}}Vw_{i}^{V}\mu_{i}^{V} + \summation_{i=1}^{n_{c}}Hw_{i}^{H}\mu_{i}^{H}
\end{displaymath}
%\end{mdframed}


Now replacing Eqn.~\ref{Eqn_MassBalance_MassFraction2} in Eqn.~\ref{MassGibbs_Definition3},
\begin{equation}
g = \summation_{i=1}^{n_{c}-1}\Pi^{1}w_{i}^{1}\left(\mu_{i}^{1}-\mu_{n_{c}}^{1}\right) + \Pi^{1}\mu_{n_{c}}^{1} + \summation_{i=1}^{n_{c}-1}\Pi^{2}w_{i}^{2}\left(\mu_{i}^{2}-\mu_{n_{c}}^{2}\right) + \Pi^{2}\mu_{n_{c}}^{2} + \cdots + \summation_{i=1}^{n_{c}-1}\Pi^{n_{p}}w_{i}^{n_{p}}\left(\mu_{i}^{n_{p}}-\mu_{n_{c}}^{n_{p}}\right) + \Pi^{n_{p}}\mu_{n_{c}}^{n_{p}}
\label{MassGibbs_Definition4}
\end{equation}
Now replacing Eqns.~\ref{Eqn_MassBalance_PhaseMassFraction2} and~\ref{Eqn_MassBalance_FeedMassFractionConstraint2} in~\ref{MassGibbs_Definition4}:
\begin{eqnarray}
g &=& \summation_{i=1}^{n_{c}}\Pi^{1}w_{i}^{1}\left(\mu_{i}^{1}-\mu_{n_{c}}^{1}\right) + \Pi^{1}\mu_{n_{c}}^{1} + \cdots + \summation_{i=1}^{n_{c}}\Pi^{n_{p}-1}w_{i}^{n_{p}-1}\left(\mu_{i}^{n_{p}-1}-\mu_{n_{c}}^{n_{p}-1}\right) + \Pi^{n_{p}-1}\mu_{n_{c}}^{n_{p}-1} + \nonumber \\
   && \summation_{i=1}^{n_{c}-1}\left[z_{i} - \summation_{j=1}^{n_{p}-1}\Pi^{j}w_{i}^{j}\right]\left(\mu_{i}^{n_{p}}-\mu_{n_{c}}^{n_{p}}\right) + \Pi^{n_{p}}\mu_{n_{c}}^{n_{p}}\nonumber
\end{eqnarray}

Rearranging
\begin{eqnarray}
g &=& \summation_{i=1}^{n_{c}}\Pi^{1}w_{i}^{1}\left[\left(\mu_{i}^{1}-\mu_{i}^{n_{p}}\right) - \left(\mu_{n_{c}}^{1}-\mu_{n_{c}}^{n_{p}}\right)\right] + \cdots + \summation_{i=1}^{n_{c}}\Pi^{n_{p}-1}w_{i}^{n_{p}-1}\left[\left(\mu_{i}^{n_{p}-1}-\mu_{i}^{n_{p}}\right) - \left(\mu_{n_{c}}^{n_{p}-1}-\mu_{n_{c}}^{n_{p}}\right)\right] + \nonumber \\
&& \summation_{j=1}^{n_{p}-1}\Pi^{j}\mu_{n_{c}}^{j} + \summation_{i=1}^{n_{c}-1}z_{i}\left(\mu_{i}^{n_{p}}-\mu_{n_{c}}^{n_{p}}\right) +\left(1-\summation_{j=1}^{n_{p}-1}\Pi^{j}\right)\mu_{n_{c}}^{n_{p}} \nonumber
\end{eqnarray}

With a compressed notation:

\begin{eqnarray}
g &=& \summation_{i=1}^{n_{c}-1}\left\{\summation_{j=1}^{n_{p}-1}\Pi^{j}w_{i}^{j}\left[\left(\mu_{i}^{j}-\mu_{i}^{n_{p}}\right) - \left(\mu_{n_{c}}^{j}-\mu_{n_{c}}^{n_{p}}\right)\right]\right\} + \summation_{j=1}^{n_{p}-1}\Pi^{j}\left(\mu_{n_{c}}^{j}-\mu_{n_{c}}^{n_{p}}\right) +  \nonumber \\
&& \hspace{3cm} \summation_{i=1}^{n_{c}-1} z_{i}\mu_{i}^{n_{p}} + \underbrace{\left(1-\summation_{i=1}^{n_{c}-1}z_{i}\right)}_{\red{z_{n_{c}}}}\mu_{n_{c}}^{n_{p}} \nonumber 
\end{eqnarray}
Thus
\red{\begin{equation}
g = \summation_{i=1}^{n_{c}-1}\left\{\summation_{j=1}^{n_{p}-1}\Pi^{j}w_{i}^{j}\left[\left(\mu_{i}^{j}-\mu_{i}^{n_{p}}\right) - \left(\mu_{n_{c}}^{j}-\mu_{n_{c}}^{n_{p}}\right)\right]\right\} + \summation_{j=1}^{n_{p}-1}\Pi^{j}\left(\mu_{n_{c}}^{j}-\mu_{n_{c}}^{n_{p}}\right) + \summation_{i=1}^{n_{c}}z_{i}\mu_{i}^{n_{p}}
\end{equation}}

The chemical potential was defined as a functional (Eqn.~\ref{chempotential_functional}):
\begin{eqnarray}
&& \mu_{i}^{1} = \mu_{i}^{1}\left(w_{1}^{1},w_{2}^{1},\cdots,w_{n_{c}}^{1}\right) \nonumber \\
&& \mu_{i}^{2} = \mu_{i}^{2}\left(w_{1}^{2},w_{2}^{2},\cdots,w_{n_{c}}^{2}\right) \nonumber \\
&& \hspace{2cm}\vdots \nonumber \\
&& \mu_{i}^{n_{p}} = \mu_{i}^{n_{p}}\left(w_{1}^{n_{p}},w_{2}^{n_{p}},\cdots,w_{n_{c}}^{n_{p}}\right) \label{ChemPotDef2}
\end{eqnarray}
with constraint Eqn.~\ref{Eqn_MassBalance_FeedMassFractionConstraint2} in Eqn.~\ref{ChemPotDef2},
\begin{eqnarray}
\mu_{l}^{k} &=& \mu_{l}^{k}\left(w_{i}^{j},\Pi^{j}\right) \nonumber \\
           &=& \mu_{l}^{k}\left(w_{1}^{1},\cdots,w_{n_{c}-1}^{1},w_{1}^{2},\cdots,w_{n_{c}-1}^{2},\cdots,w_{1}^{n_{p}},w_{2}^{n_{p}},\cdots,w_{n_{c}-1}^{n_{p}},\Pi^{1},\cdots,\Pi^{n_{p}}\right)\nonumber
\end{eqnarray}
with $i = 1, 2, \cdots, l, \cdots, n_{c} - 1$, $j = 1, 2, \cdots, n_{p}$ and $j\neq k$

And the mass-based free Gibbs energy should take the following format
\begin{displaymath}
g = g\left(w_{i}^{j},\Pi^{j}\right).
\end{displaymath}
From classical thermodynamic, $\mu_{i}^{j}$ can be defined as a function of the fugacity coefficient, $\varphi_{i}^{j}$,
\begin{equation}
\mu_{i}^{j}= R T\left[\ln\varphi_{i}^{j} + \ln\left(P x_{i}^{j}\right)\right] + \theta_{i}\left(T\right),
\end{equation}
where $\theta_{i}$ is an integration constant. The fugacity coefficient is defined as,
\begin{equation}
\varphi_{i}^{j} = \frc{f_{i}^{j}}{P x_{i}^{j}},
\end{equation}
where $f$ and $x$ are the fugacity and molar fraction respectively.


 % Compositional

%
%%%
%%% CHAPTER
%%%
\chapter{Hydrate Formation}\label{Chapter:Hydrate}


%%%
%%% Section
%%%
\section{Mass Conservation and Degrees of Freedom Analysis}\label{Chapter:Hydrate:Section:MassConservation}

Assuming an isothermal flash with the following phases present in the equilibrium:
  \begin{center}
    \begin{tabular}{l l l l}
      $L$: & liquid phase & $H$: & hydrate phase \\
      $V$: & vapour phase & $I$: & ice phase \\ 
    \end{tabular}
  \end{center}
$L$, $V$, $H$ and $I$ are the molar fraction of each phase. In addition, $F$ represents the molar fraction of the feeding stream. Finally, let's define $x_{i}^{j}$ as the molar fraction of component $i \left(=1,2,\cdots,n_{c}\right)$ in phase $j \left(=L,V,H,I,F\right)$. Thus, assuming a closed system (\ie no mass transfer across a defined domain border):
  \begin{eqnarray}
     && L + V + H + I = 1 \\\label{Chapter:Hydrate:Eqn:Balance:MassFractionPhase}
     && \mfr[x]{i}{F}F = \mfr[x]{i}{V}V + \mfr[x]{i}{L}L + \mfr[x]{i}{H}H + \mfr[x]{i}{I}I, \label{Chapter:Hydrate:Eqn:Balance:MassFractionComponent}
  \end{eqnarray}
with the following linear constraints:
  \begin{equation}
    \summation_{i}^{n_{c}} \mfr[x]{i}{j} = 1, \hspace{2cm} j=L,V,H,I,F\label{Chapter:Hydrate:Eqn:Balance:MassFractionConstraint}
  \end{equation}
 
The resulting system of non-linear algebraic equations have $4n_{c}+6$ unknowns, \ie $\mfr[x]{i}{V}, \mfr[x]{i}{L}, \mfr[x]{i}{H}, \mfr[x]{i}{I}$ (each of them with $n_{c}$ components), $L$, $V$, $I$, $H$, temperature ($T$) and pressure ($P$).  In the case of methane hydrate, $n_{c}=2$ $\left(\text{\ie} CH_{4}+H_{2}O\right)$. Additionally, the ice phase ($I$) contains just water -- $w_{w}^{I}=1$, and the number of unknowns and equations (linear constraint for ice) are reduced.

In the equilibrium, the equality of phase fugacities for a given component $i$,
   \begin{equation}
      \mfr[f]{i}{V} = \mfr[f]{i}{L} = \mfr[f]{i}{H} = \mfr[f]{i}{I} \hspace{2cm}\text{with } i=1,\cdots,n_{c} \label{Chapter:Hydrate:Eqn:FugacityEquality}
   \end{equation}
Note that one of the relations in Eqn.~\ref{Chapter:Hydrate:Eqn:FugacityEquality} is not independent. Thus the total number of independent equations is $4n_{c}+4$, and the number of degrees of freedom $\left(\varphi\right)$, \ie   
   \begin{center} 
      (number of unknowns) - (number of independent equations),
   \end{center}
is equal to $2$.

If we specify the pair $\left[T,P\right]$, we can determine the composition of each component, $\mfr[x]{i}{j}$, and each phase $j$. The Gibb's phase rule,
  \begin{eqnarray}
     \varphi &=& 2 + n_{c} - n_{p} \\ \nonumber
             &=& 2 + 2 - n_{p} = 4 - n_{p},
  \end{eqnarray}
states the possibility of a quadruple point $\left(\varphi=0\right)$, however if we want to specify $\left[T,P\right]$ then a two-phase system may occur.



%%%
%%% Section
%%%
\section{Fugacity Models}\label{Chapter:Hydrate:Section:FugacityModels}

%%%
\subsection{Models for Liquid and Vapour Phases}\label{Chapter:Hydrate:Section:FugacityModels:Section:L_V}

%%%
\subsubsection{Soave-Redlich-Kwong Equation of State (SRK-EOS)}\label{Chapter:Hydrate:Section:FugacityModels:Section:L_V:Section:SRKEOS}
The cubic polynomial in $z$ (compressibility factor) representation of the SRK-EOS can be expressed as,
   \begin{equation}
      z^{3} - z^{2} + \left(A-B-B^{2}\right)z - AB = 0 \label{Chapter:Hydrate:Section:FugacityModels:Section:L_V:Section:SRKEOS:Eqn:SRKEOS}
   \end{equation}
where $A$ and $B$ are defined in Eqn.~\ref{Chapter:EOSPR:Section:PR:Eqn:CubicPR_Z_Parameters} with
   \begin{displaymath}
      a = \sum\limits_{i=1}^{n_{c}}\sum\limits_{j=1}^{n_{c}}x_{i}x_{j}a_{i,j} \hspace{1cm}\text{ and }\hspace{1cm} b = \sum\limits_{i=1}^{n_{c}}x_{i}b_{i}
   \end{displaymath}
with
   \begin{displaymath}
      a_{i}=0.42747\frc{\left(RT_{c,i}\right)^{2}}{P_{c,i}}, \hspace{1cm} b_{i}=0.08664\frc{RT_{c,i}}{P_{c,i}}\hspace{1cm}\text{ and }\hspace{1cm}a_{i,j} = \sqrt{\alpha_{i}a_{i}\alpha_{j}a_{j}}
   \end{displaymath}

The fugacity is  expressed as 
   \begin{eqnarray}
      f_{i}(T,P) &=& P\exp{\left[\left(z-1\right)-\ln{\left(z-B\right)}-\frc{A}{B}\ln{\left(1+\frc{B}{Z}\right)}\right]}  \\
      f\left(T,P,x_{i}\right) &=& x_{i}P\exp{\left[\frc{b_{i}}{b}\left(z-1\right)-\ln{\left(z-B\right)}-\frc{A}{B}\left(\frc{2\sum\limits_{j=1}^{n_{c}}x_{j}a_{i,j}}{a}-\frc{b_{i}}{b}\right)\ln{\left(1+\frc{B}{Z}\right)}\right]} 
   \end{eqnarray}

%%%
\subsubsection{Peng-Robinson Equation of State (PR-EOS)}\label{Chapter:Hydrate:Section:FugacityModels:Section:L_V:Section:PREOS}
The cubic polynomial in $z$ (compressibility factor) representation of the SRK-EOS can be expressed as,
   \begin{displaymath}
      z^{3} -\left(1-B\right)z^{2} + \left(A-2B-3B^{2}\right)z^{2} - \left(AB -B^{3} - B^{2}\right) = 0
   \end{displaymath}
(see Eqn.~\ref{Chapter:EOSPR:Section:PR:Eqn:CubicPR_Z}) with fugacities,
   \begin{eqnarray}
      f_{i}(T,P) &=& P\exp{\left\{\left(z-1\right)-\ln{\left(z-B\right)}-\frc{A}{2\sqrt{2}B}\ln{\left[\frc{z + \left(\sqrt{2}+1\right)}{z-\left(\sqrt{2}-1\right)}\right]}\right\}}  \\
      f\left(T,P,x_{i}\right) &=& x_{i}P\exp\left\{\frc{b_{i}}{b}\left(z-1\right)-\ln{\left(z-B\right)}-\frc{A}{1\sqrt{2}B}\left(\frc{2\sum\limits_{j=1}^{n_{c}}x_{j}a_{i,j}}{a}-\frc{b_{i}}{b}\right) \times \right.\nonumber \\
                              && \hspace{6cm}\left.\ln{\left[\frc{z + \left(1-\sqrt{2}\right)B}{z +\left(1+\sqrt{2}\right)B}\right]}\right\} 
   \end{eqnarray}

%%%
\subsection{Host-Component in the Liquid Phase}\label{Chapter:Hydrate:Section:FugacityModels:Section:HostLiquid}
The host component (hydrocarbon) in the liquid phase is constant for a given temperature according to the Henry's law for the solubility calculation of hydrocarbons in water -- negligible at low temperatures bur relatively high for high temperatures,
   \begin{equation}
      \mfr[x]{i}{L} = \mfr[f]{i}{L}\left[\mathcal{H}_{i}\exp{\left(\frc{P\overline{V}^{\infty}}{RT}\right)}\right]^{-1},\hspace{1cm} i=1,2,\cdots,n_{c}
   \end{equation}
where $\overline{V}^{\infty}$ is the partial molar volume at infinite dilution, $\mathcal{H}_{i}$ is the Henry's constant~\citep{Prausnitz_2000} defined as
   \begin{equation}
      -\ln{\left[\frc{\mathcal{H}_{i}}{101325}\right]} = \mfr[\mathcal{H}]{i}{1} + \frc{\mfr[\mathcal{H}]{i}{2}}{T} + \mfr[\mathcal{H}]{i}{3}\ln{T} +  \mfr[\mathcal{H}]{i}{4}T,
   \end{equation}
where $\mfr[\mathcal{H}]{i}{j}$ are constants. For binary systems, the molar fraction of water in the liquid phase can be calculated as,
   \begin{equation}
     \mfr[x]{w}{L} = 1 - \mfr[f]{i}{L}\left[\mathcal{H}_{i}\exp{\left(\frc{P\overline{V}^{\infty}}{RT}\right)}\right].
   \end{equation}




%%%
%%% Section
%%%
\section{Models for the Hydrate Phase}\label{Chapter:Hydrate:Section:HydrateModels}


%%%
\subsection{van der Waals and Platteeuw Model}\label{Chapter:Hydrate:Section:HydrateModels:Section:vdW_Platteeuw}
This model~\citep{vdWP_1959} is based on the assumption of the equality of the chemical potential of water in the hydrate phase ($H$) and a second phase, either liquid ($L$) or ice ($I$). We can generalise (for simplicity in the notation) these two phases as $\Pi$:
   \begin{equation}
      \Delta\mfr[\mu]{w}{H} = \mfr[\mu]{w}{\beta} - \mfr[\mu]{w}{H} =  \mfr[\mu]{w}{\beta} - \mfr[\mu]{w}{\Pi} = \Delta\mfr[\mu]{w}{\Pi},
   \end{equation}
where $\mfr[\mu]{w}{\beta}$ is the chemical potential of water when the clathrate cage is {\it empty} -- this is often considered as a meta-stable phase as the hydrate structure can still `welcome' a guest-molecule (hydrocarbon). And,
   \begin{equation}
     \mfr[\mu]{w}{H}(T,P) = \mfr[\mu]{w}{\beta}(T,P) + KT\summation_{m}\nu_{m}\ln{\left[1-\summation_{i}\Theta_{im}(T,P)\right]}
   \end{equation}

%%%
\subsection{Fugacity of Water}\label{Chapter:Hydrate:Section:HydrateModels:Section:FugacityWater}
\citet{klauda_2000}:
   \begin{equation}
      \mfr[f]{w}{H}(T,P) = \mfr[f]{w}{\beta}(T,P)\exp{\left[-\frc{\Delta\mfr[\mu]{w}{H}(T,P)}{RT}\right]},
   \end{equation}
with 
   \begin{eqnarray}
       \Delta\mfr[\mu]{w}{H} = -RT\summation_{m}\nu_{}\ln{\left(1-\summation_{i}\Theta_{im}\right)}\hspace{2cm}\text{ and}, \\
       \mfr[f]{w}{\beta}(T,P) = \mfr[P]{w}{\text{sat},\beta}(T)\exp{\left\{\frc{\mfr[v]{w}{\beta}\left[P-\mfr[P]{w}{\text{sat},\beta}(T)\right]}{RT}\right\}},
   \end{eqnarray}
where $\mfr[P]{w}{\text{sat},\beta}$ is the vapour pressure of the water in the hydrate phase, $\mfr[v]{w}{\beta}$ is the volume of the empty hydrate that depends on the hydrate structure, $k$ is the Boltzmann constant, $\nu_{m}$ is the number of cavities (\ie cages) od the hydrate structure of type $m$ per water molecule, $\Theta_{im}(T,P)$ is the volume fraction occupied by the guest molecule $i$ in cavity $m$ and is calculated from the Langmuir adsorption isothermal relation,
   \begin{equation}
      \Theta_{im}(t,P) = \frc{C_{im}(T)f_{i}(T,P)}{1=\summation_{m}C_{m}(T)f_{i}(T,P)}
   \end{equation}
 % Simple Hydrate Model

\pagebreak

\cleardoublepage

%\part{Appendices}

%  \begin{appendix}
%     \input{Appendix_ConversionTable}
%     \input{Appendix_SaturatedTables}
%     \input{Appendix_LogTable}
%     \input{Appendix_Calculus1}
%     \input{Appendix_NumMethods}
%     \input{Appendix_Examples}
%  \end{appendix}

\cleardoublepage

\pagebreak

%\bibliographystyle{plain}
\bibliography{references}
\bibliographystyle{unsrt}

\cleardoublepage
\phantomsection
\renewcommand\leftmark{}
\renewcommand\rightmark{Index}
\addcontentsline{toc}{chapter}{Index}

\printindex


\end{document}
